\section{The full-addition base product and its surviving support geometry}
\label{sec:full-support-base-product}

The supplied arithmetic-specificity calculation computes the self-product
of the scalar globalization. Here we retain the entire finite lattice
of the original version 11 definition \cite{support}. The result specifies
which coordinates agree in the product, which remain independent, and
the new relative cohomology they support. The blueprint conventions are
those of Lorscheid \cite{lorscheid}, original author TeX
\path{blueprints-v2.tex}, axioms (A1)--(A9),
\texttt{lemma\_proper}, and Section 1.12's tensor construction.
We use the term properization for its quotient by equality relations
between individual elements. Semiring completion also adds the values
of formal sums; the two operations below have different explicit images.

\subsection{The integer coordinate in every common target}

Put $S=G(\mathbb Z)$ and denote its supported integers by $j(n)$,
with $j(0)=e$ and its additive identity $\tau$. For maps $f,g$ from
the full-addition blueprint of $S$ to a common proper blueprint, set
$u=f(j(-1))$ and $v=g(j(-1))$. Both satisfy the original relations
\[
 1+u+1\equiv1,\qquad u+1+u\equiv u,
 \qquad 1+v+1\equiv1,\qquad v+1+v\equiv v.
 \tag{BP1}
\]
In the commutative monoid of formal sums modulo preaddition,
these regular inverses of $1$ are unique. Indeed
\[
 u\equiv u+1+v+1+u\equiv u+1+v,
 \qquad v\equiv v+1+u+1+v\equiv u+1+v.
 \tag{BP2}
\]
The first equality in each line replaces $1$ by $1+v+1$ or
$1+u+1$ inside the corresponding regular-inverse identity;
the second uses $u+1+u\equiv u$ or $v+1+v\equiv v$.
Properness implies $u=v$. Then $f(e)\equiv1+u\equiv g(e)$,
so these also agree. Repeated addition determines every positive
integer from $1$ and every negative integer from $u$; both maps send
$\tau$ to the unique zero. Thus $f=g$ on the entire scalar $S$.
This proof uses the regular inverse, not an additive inverse of $1$
with respect to $\tau$.

Consequently, for any specified blueprint base $B\to S$,
\[
 (B_S\otimes_B B_S)^{\rm prop}\simeq B_S,
 \qquad S\otimes_{\mathbb N}S\simeq S.
 \tag{BP3}
\]
For either assertion, pairs of agreeing maps into the relevant common
target are exactly single maps, by (BP2). The inverse arrows are
the inclusion $a\mapsto a\otimes1$ and multiplication
$a\otimes b\mapsto ab$. Their composites are identities by that
universal property. This proves the scalar calculation in the arrival.
In particular $e\otimes1=1\otimes e$; the actual prime $(e)$ remains
present and does not become two different arithmetic coordinates.

\subsection{The full lattice product in commutative semirings}

Let $P,Q$ be nonempty finite posets, with their topology of lower
sets. Put $L=\operatorname{Down}(P)$, $M=\operatorname{Down}(Q)$,
and $T=\operatorname{Down}(P\times Q)$. The exact coproduct is
\[
 G_L(\mathbb Z)\otimes_{\mathbb N}G_M(\mathbb Z)
       \simeq G_T(\mathbb Z).
 \tag{BP4}
\]
The structure maps send every supported integer to the same integer
and send the zero-amplitude supports by
\[
 A\longmapsto A\times Q,\qquad B\longmapsto P\times B.
 \tag{BP5}
\]
Their product on zero-amplitude elements has support $A\times B$.
All formulas preserve both $\tau$ and supported $e$ in their exact
roles.

Here is a proof of the whole universal property. Given semiring maps
$f:G_L(\mathbb Z)\to C$ and $g:G_M(\mathbb Z)\to C$, (BP2)
makes their supported integer maps agree. Write $E$ for their common
value of $e$, and $a_A=f(0,A)$, $b_B=g(0,B)$. Each is additively
and multiplicatively idempotent, is absorbed by $E$ under addition,
and has $E$ as a multiplicative identity. The finite sums of their
products form a bounded distributive lattice with join given by
addition, meet by multiplication, bottom $0_C$ and top $E$.
For detail, sums of commuting multiplicative idempotents which are
also additively idempotent remain idempotent under either operation.
Cross absorption holds since $b_B+E=E$ implies
$a_Ab_B+a_AE=a_AE=a_A$. Distributivity and this identity prove both
absorption laws on the generated set. The two original lattice maps
therefore take values in this common bounded lattice.

The universal lattice map from $T$ has the explicit value
\[
 F(U)=\bigvee_{(p,q)\in U} a_{\downarrow p}\,b_{\downarrow q}.
 \tag{BP6}
\]
It preserves unions immediately. For intersections, distribute the
two joins and use
$\downarrow p\cap\downarrow p'
 =\bigcup_{r\le p,p'}\downarrow r$, and the identical formula
in $Q$. The resulting terms are exactly those indexed by $U\cap U'$.
It preserves bottom and top. The same expansion proves
$F(A\times Q)=a_A$ and $F(P\times B)=b_B$. Every lower set is the
union of its principal rectangles, so (BP6) is the unique such map.

Extend $F$ to $G_T(\mathbb Z)$ by the common supported integer map.
A supported integer absorbs each generator under addition and
acts identically on it by multiplication; these identities hold
first in its original factor and then for products and joins.
They also hold for a negative supported integer, because both factors
have the identical map $j(-1)$ proved in (BP2). Thus the extension
preserves the mixed operations. It agrees at $e$, so is well-defined.
It is the required unique semiring map. This proves (BP4).

For arbitrary finite nontrivial bounded distributive lattices this
description includes every case, rather than a chosen subfamily.
To see this directly, write $P$ for the nonzero join-irreducible
elements of $L$, with inherited order. Every element is the join
of the join-irreducibles below it, by induction in the finite order.
Distributivity implies a join-irreducible below a finite join lies
below one summand: for two summands intersect it with the join and
use its defining irreducibility, then induct. Hence
$a\mapsto\{p:p\le a\}$ is a lattice isomorphism
$L\simeq\operatorname{Down}(P)$, with inverse the join.
The same proof applies to $M$.

\subsection{Proper blueprint product, semiring completion, and their exact map}

Use the base blueprint with zero $\mathbb F_1^{\rm mon}=\{0,1\}$,
with only its zero relation, and let $\mathcal B$ be the raw
blueprint tensor of the two full-addition factors. In the description
of (BP4), its properization is the subblueprint with underlying set
\[
 \mathcal A_{P,Q}
 =\{(n,P\times Q):n\in\mathbb Z\}
       \ \cup\ \{(0,A\times B):A\in L,\ B\in M\},
 \tag{BP7}
\]
and preaddition the equality of its finite sums in $G_T(\mathbb Z)$.
In particular the exact connecting maps are
\[
 \mathcal B\longrightarrow\mathcal B^{\rm prop}
   \simeq\mathcal A_{P,Q}\hookrightarrow B_{G_T(\mathbb Z)},
 \qquad (\mathcal B^{\rm prop})_{\mathbb N}=G_T(\mathbb Z).
 \tag{BP8}
\]

Indeed the raw underlying elements are pure tensors; the finite sums
modulo their preaddition form its universal semiring completion.
Maps from that completion to $C$ are the two original semiring maps,
so (BP4) identifies this completion exactly. A pure tensor has the
value $(nm,P\times Q)$ when both amplitudes are nonzero, and otherwise
has amplitude zero and rectangular support as in (BP7). All elements
of (BP7) occur. Equality of images of two raw elements in the semiring
completion is exactly their preaddition equivalence, by its definition
as the quotient of formal sums. Properization consequently has
precisely (BP7)'s elements and induced preaddition. Adding all their
finite sums gives all lower sets by principal rectangles, proving
the last assertion of (BP8). Multiplication of rectangles closes
(BP7), so its indicated subblueprint is defined.

This identifies a material extra operation. If $P=Q$ consists of
two incomparable points, the union
$\{(p_1,p_1),(p_2,p_2)\}$ is a lower set but is not a rectangle.
Its zero-amplitude value is a sum in the semiring completion, while
it is absent from the underlying set (BP7). This is an explicit
example and exact embedding, with no identification of properization
with adding all sums. For scalar one-point $P,Q$, every lower set
is a rectangle and (BP8) reduces to (BP3).

\subsection{The changed prime spectrum and the diagonal}

The full spectrum theorem already proved in this manuscript gives
\[
 \operatorname{Spec}G_T(\mathbb Z)
     =(P\times Q)\mathbin{\blacktriangleright}\operatorname{Spec}\mathbb Z.
 \tag{BP9}
\]
For completeness the lattice prime associated to $(p,q)$ is
$\{U:(p,q)\notin U\}$: a union avoids the point exactly when
both sets do, and an intersection avoids it exactly when one does.
Every finite lattice prime has this form, since its complement is
a prime filter, whose smallest element is a principal lower set;
the join-irreducibility proof following (BP6) identifies its generator.
The general support/arithmetic classification then proves (BP9).
The two projection maps on spectra are
\[
 (p,q)\mapsto p,\quad (p,q)\mapsto q;
 \qquad \mathfrak p\mapsto\mathfrak p
       \quad(\mathfrak p\in\operatorname{Spec}\mathbb Z).
 \tag{BP10}
\]
Testing inverse images against (BP5) proves these assertions point
by point. There are no support/arithmetic mixed points: the two
pulled-back supported zeros are the same element $e$, so any prime
contains both or neither. Every integer and its original prime
vanishing set is retained on the single arithmetic stratum.

For $P=Q$, multiplication of the two equal factors induces the
surjective map
\[
 G_{\operatorname{Down}(P\times P)}(\mathbb Z)
       \longrightarrow G_{\operatorname{Down}(P)}(\mathbb Z),
 \qquad (0,U)\mapsto(0,\{p:(p,p)\in U\}),
 \quad j(n)\mapsto j(n).
 \tag{BP11}
\]
It preserves joins and meets because it is inverse image under
the diagonal, so is a semiring map; it is surjective using $A\times P$.
Its spectrum map is the diagonal on the full support and the identity
on the arithmetic stratum. Its exact congruence identifies two null
supports precisely when their diagonal intersections agree, and
identifies no distinct nonzero amplitudes. This gives the complete
defect object and comparison map of the self-product calculation.

\subsection{An explicit product with surviving relative cohomology}

Take $P$ to be the six-element face poset of the boundary of an
oriented triangle: three vertices and its three edges, ordered by
incidence. This is the earlier (FL) support example. The poset
$P\times P$ is the face poset of the product cell structure on
the product of the two boundary triangles. Its order complex is
the barycentric subdivision of that product. An explicit realization
sends a pair of faces to the product of their barycentres and extends
affinely on each chain. On each product face, the chains ending in
that face form the cone from its barycentre over the subdivided
boundary. Induction over dimensions $0,1,2$ shows these cones tile
each convex product face with disjoint interiors and agree on common
boundaries. Thus this realization is a homeomorphism, including every
one of the $36$ original support points.

Here is also a direct integral calculation. The triangle boundary
has cellular boundary $\partial e_i=v_{i+1}-v_i$ (indices modulo
three). Choose $e_0,e_1$ as a spanning tree. Their boundaries and
$v_0$ form an integral basis of the vertex group; the edge basis
$e_0,e_1,z=e_0+e_1+e_2$ splits the chain complex into two identity
complexes and a copy of $\mathbb Z$ in each degree $0,1$.
The product cells have boundary
$\partial(a\times b)=\partial a\times b+(-1)^{\deg a}a\times\partial b$,
as follows by orienting each rectangular face by its two edge
coordinates. Tensoring the displayed splitting leaves its two
identity factors contractible and gives
\[
 H^j(P\times P,\mathbb Q)=
 \begin{cases}\mathbb Q&j=0,2,\\\mathbb Q^2&j=1,\\0&j>2.\end{cases}
 \tag{BP12}
\]
Write $\alpha_1,\alpha_2$ for the edge-circuit cohomology classes
pulled back from the factors, each evaluating to one on $z$.
Their product evaluates to one on the oriented product $z\times z$.
The diagonal pulls both to $\alpha$ and their degree-two product to
zero, since the one-dimensional boundary has no degree-two cochains.

Apply the actual relative complex (FL14),
$[\mathbb Z\to C^\bullet(P\times P,\mathbb Q)]$, with the
integer term in degree zero and the support terms shifted by one.
The constants map is injective, and (BP12) gives
\[
 H_X^1=\mathbb Q/\mathbb Z,\qquad
 H_X^2=\mathbb Q\alpha_1\oplus\mathbb Q\alpha_2,\qquad
 H_X^3=\mathbb Q(\alpha_1\alpha_2),\qquad H_X^{j>3}=0.
 \tag{BP13}
\]
The morphism (BP11) pulls degree one identically, degree two by
$(a,b)\mapsto a+b$, and degree three to zero. Thus the support
product has two independent degree-two relative classes and one
degree-three class, and the diagonal's exact lost classes are
$\mathbb Q(\alpha_1-\alpha_2)$ and $\mathbb Q(\alpha_1\alpha_2)$.
They are produced by the complete original support product. Their
existence establishes these groups and maps; a Frobenius action or
a positive Weil pairing is not assigned to them without a construction.

\begin{figure}[ht]
\centering
\begin{tikzcd}[column sep=large,row sep=large]
(P\times P)\blacktriangleright\operatorname{Spec}\mathbb Z
  \arrow[r,shift left,"\mathrm{pr}_1"]
  \arrow[r,shift right,swap,"\mathrm{pr}_2"]
 &P\blacktriangleright\operatorname{Spec}\mathbb Z
  \arrow[l,bend left=35,"\mathrm{diag}"]\\
\mathbb Q\alpha_1\oplus\mathbb Q\alpha_2
  \arrow[r,"{(a,b)\mapsto a+b}"]
 &\mathbb Q\alpha
\end{tikzcd}
\caption{The full-support product and its diagonal, (BP9)--(BP13).
The top row displays maps of spectra; the bottom displays the pullback
of relative degree-two cohomology along the diagonal. Both arithmetic
maps are the identity. For the six-point triangle-boundary support,
the product retains all $36$ support points and the additional relative
degree-three class. Sources of the coefficient object and blueprint
operations: version 11 \cite{support} and Lorscheid \cite{lorscheid}.}
\end{figure}
